\begin{table}[htbp]\centering\small
\caption{Ingresos presupuestarios: monto y razón a PIB, línea G}\label{cua:ingpib}
\begin{tabular}{lrrrr}
\toprule
Concepto & 2026 aprob. (mdp) & 2027 (mdp) & 2026 (\% PIB) & 2027 (\% PIB) \\
\midrule
Ingresos presupuestarios & 8.721.057,3 & 9.156.528,9 & 22,5 & 23,2 \\
\quad Petroleros & 1.204.277,7 & 984.539,7 & 3,1 & 2,5 \\
\quad No petroleros & 7.516.779,6 & 8.171.989,2 & 19,4 & 20,7 \\
\quad\quad Tributarios & 5.838.571,0 & 6.263.886,8 & 15,1 & 15,9 \\
\quad\quad No tributarios & 377.130,1 & 521.763,5 & 1,0 & 1,3 \\
\quad\quad Organismos y empresas & 1.301.078,5 & 1.386.338,9 & 3,4 & 3,5 \\
\bottomrule
\end{tabular}
\\[2pt]\begin{minipage}{0.92\textwidth}\scriptsize Fuente: CGPE 2027, edición SHCP, Anexo II.6 (p. 67). Las razones de 2026 usan el PIB \emph{aprobado} de 2026 (38.715,9 mmp), como declara la nota del propio anexo, y no el estimado.\end{minipage}
\end{table}
